%  $Author: awl8049 $$ 
%  $Date: 2011/12/13 17:51:16 $
%  $Revision: 3.5 $
% 
\documentclass[10pt]{letter} % Uses 10pt
%Use \documentstyle[newcent]{letter} for New Century Schoolbook postscript font
% the following commands control the margins:
\usepackage[paper=letterpaper,
            marginparwidth=0in,
            marginparsep=0.5in,
            margin=1in,
            includemp]{geometry}

% Add the code to give us a biblography environment
\usepackage{enhletter}
% Comment environment support
\usepackage{comment}
\usepackage{graphicx}
% Use fancyhdr and lastpage to get decent page headers and footers
\usepackage{lastpage}
\usepackage{fancyhdr}
\pagestyle{fancy}
\fancyhf{}
\lhead{}
\chead{}
\rhead{}
\lfoot{TACO-2010-27: Review Response: Reviewer \#\thesection, Page \thepage}
\cfoot{}
\rfoot{}
\renewcommand{\headrulewidth}{0pt}
\renewcommand{\footrulewidth}{2pt}

\newenvironment{rviewcomment}
{~\\%
\begin{bfseries}}
{\end{bfseries}}

\newcommand{\rviewresponse}{\textbf{\textit{Response}}:}
\newcommand{\rviewresponses}{\textbf{\textit{Responses}}:}

% Packages for managing per section biblographies
\usepackage[sectionbib]{chapterbib}

\begin{document}

%\signature{Adam Wade Lewis}           % name for signature 
\longindentation=0pt                       % needed to get closing flush left
\let\raggedleft\raggedright                % needed to get date flush left
 
 
\begin{letter}{Professor Tom Conte \\
Editor-in-Chief \\
ACM Transactions on Architecture and Code Optimization }


\begin{flushright}
{\hfill \large\bf The University of Louisiana at Lafayette}
\end{flushright}
\medskip\hrule height 1pt
\begin{flushright}
\hfill Center for Advanced Computer Studies\\
\hfill P.O. Box 44330 \\
\hfill Lafayette, Louisiana 70504-4330\\
\hfill (337) 482-6338 \\
\hfill \email{awlewis@cacs.louisiana.edu}
\end{flushright} 
%\vfill % forces letterhead to top of page

 
\opening{Dear Professor Conte:} 
 
\noindent We thank the reviewers for the time and effort required for
their thorough reviews of our revision and greatly appreciate their
candid and valuable comments, which clearly strengthen our manuscript.
We have taken all comments into full consideration, addressing them in
the revised manuscript.  Please find a detailed
response to the review comments raised by Reviewer \#1, in the enclosed
attachment of this cover letter.
 
\closing{Respectfully yours, \\
\fromsig{\includegraphics[scale=0.2]{sig.jpg}} \\
\fromname{Adam Wade Lewis}}

\encl{(1) Response to Reviewer \#1}
\clearpage
\section{TACO-2010-27: Review Response: Reviewer \#1}
\label{sec-1}
First of all, we greatly appreciate your time and effort put toward
your candid and informative review, which helps to improve the quality
of our submission.  The following lists our responses to your concerns
and comments made on our earlier revision.  Hopefully, we have addressed
all your concerns and comments to your satisfaction.

\begin{rviewcomment}
  First Comment/Concern:
\end{rviewcomment}
\begin{quote}
\begin{itshape} 
  First, again, I think the paper title needs to be clearer. As it
  stands, it refers to an "approximation of an estimation". Based on
  your response, CAP is an instantaneous power estimation method, and I
  think "Run-time Energy Consumption Estimation based on time-series
  approximation.." might be more representative.
\end{itshape}
\end{quote}

\rviewresponse

The title of our paper has been adjusted according to your suggestion.

\begin{rviewcomment}
  Second Comment/Concern:
\end{rviewcomment}
\begin{quote}
  \begin{itshape}
    Next, it is still not clear to me where the benefit of CAP comes
    from. For example, can you give an example to where simple,
    Mantis-like approaches do not perform well for the workloads you
    evaluate? 
  \end{itshape}
\end{quote}
\rviewresponse 

The issue of workloads inadequately addressed by Mantis-like approaches
has been unveiled in recent work, particularly by
\cite{Varsampoulos2010},\cite{Kansal2010}, \cite{Hsu2011}, and
\cite{McCullough2011}. In \cite{Hsu2011}, the authors examined 177
measured results of the SPECpower\_ssj2008 benchmark published from
December 2007 to August 2010 to statistically analyze the shape of the
power curves over time and evaluate the effect of the aggressive
power-management schemes introduced since the publication of
\cite{Economou2006}. Both our work and \cite{Varsampoulos2010} found
that the resulting power curves were neither linear nor convex, which
invalidate an underlying assumption behind models such as MANTIS.  For
example, in \cite{McCullough2011}, the authors found that MANTIS
(adjusted for evolution of processor performance counters since the
publication of \cite{Economou2006}) exhibited the mean prediction error of
10-14\% for the SPEC CPU2006 suite, the PARSEC multicore benchmark
suite, and synthetic micro-architecture benchmarks that focused on
components of system energy consumption.  In particular, \cite{McCullough2011}
found that MANTIS suffered from pronounced error behavior in cases of
high utilization and low IPC (for instance, the \texttt{canneal}
benchmark from the PARSEC suite, the \texttt{Bonnie} I/O benchmark and a
synthetic system stress benchmark).

The causes of this error behavior were described and discussed by recent
work cited in our paper. In particular, \cite{Kansal2010} observed that
the error behavior of SPEC CPU2006 benchmarks in both physical and
virtual machine environments is rarely uniform in distributing tasks
across multi-core processors.  It was postulated in
\cite{McCullough2011} that multi-core processors have evolved to the
point where features cannot be abstracted easily to permit linear models
to accurately predict the behavior of these systems (as attempted by
MANTIS-like modeling).

We have added additional materials to our discussion of prior work
that considers the issues raised by these references.
\newpage
\begin{rviewcomment}
  Third Comment/Concern
\end{rviewcomment}

\begin{itshape}
  As you argue in the beginning of the paper, these approaches relate
  usage information to the power of the entire system rather than its
  individual components. So what kind of improvement do we see with the
  CAP approach with subcomponent prediction?
\end{itshape}

\rviewresponse 

CAP is a full-system model, aiming to predict the individual
sub-components energy consumption for attributing these components to
the total energy consumption.  In doing so, we address the concerns
raised by \cite{Kansal2010} and \cite{McCullough2011} in regard to the
factors that contribute to non-linearity in the power curves of
multi-core processors.  As shown in Tables IX and X in our revision,
both average and maximum error for CAP are fare less as compared to
those of linear AR, MARS, and EWMA predictors, in particular when we
consider a more recent processor such as the Intel processor used in our
study.

\begin{rviewcomment}
  Fourth Comment/Concern:
\end{rviewcomment}
\begin{quote}
  \begin{itshape}
    Along the same lines of comparative evaluation, as I had mentioned
    in the original review, there are other prior approaches that
    predict future behavior based patterns/statistics. How does CAP
    compare to even simplistic approaches for future behavior prediction
    like last-value, or simple exponentially-weighted moving averaging?
    I suspect the resulting power curves would look very similar, albeit
    slightly shifted versions of what is depicted in Figures 7/8. A
    simple error analysis can demonstrate the value of CAP, in addition
    to those included for AR and MARS.
  \end{itshape}
\end{quote}
\rviewresponse

Your intuition about the shape of the power curves is correct; some
similarity exists between the shape of CAP and the power curves
predicted by exponentially-weighted moving averages (EWMA) predicted
power curves.  However, EWMA demonstrates error amounts similar to those
of AR(1) and MARS predictors, with average errors between 1.0\% and
1.8\% and maximum errors between 6.9\% and 9.2\% for the AMD Opteron
processor (Intel Nehalem: average error: 1.8\% - 5.0\%; maximum error:
27.3\% - 32.4\%). CAP differs from methods such as EMWA by using points
on the attractor to approximate the next entry in the series.  By doing
so, we lower the approximation error as compared to the error amounts of
EWMA and others.  In addition
to the extra information included in Tables IX and X, we provide
additional discussion in Section 5.3 and the Appendix on this topic.
\newpage
\begin{rviewcomment}
  Fifth Comment/Concern
\end{rviewcomment}
\begin{quote}
  \begin{itshape}
    I am also still not sure how the four SPEC benchmarks can help
    verify the sub-component power models beyond memory per your
    response. Do you see these benchmarks exercising these components at
    different rates and creating different dynamic power? Section 5.1
    para-1, suggests these four were selected as workloads more common
    to server workloads. Can you please explain what this means? Can you
    also further provide an example to sub-component power estimations?
  \end{itshape}
\end{quote}
\rviewresponse

Benchmark selection in our work is based upon the decision criteria
listed in \cite{Phansalkar2007}, with the additional criterion of
selecting benchmarks from the SPEC CPU2006 suite most representative of
scientific computing workloads \cite{Cisco2010}.  These benchmarks, as
you observed, primarily focus on the processor, cache, and memory;
however, the execution of them affects the entire system, as noted in
\cite{Ye2007} and confirmed by examining the details of the PeCs
collected in our study.  The following table contains averages of the
fan speed, the number of HyperTransport bus transactions, and the amount
of bytes read/written to disks for the four benchmarks over the same
period during their execution on the AMD Operton system in our testbed:

% BEGIN RECEIVE ORGTBL pecs
\begin{tabular}{lrrrrr}
\hline
Benchmark & Fan Speed & HT0 & HT1 & HT2 & Disk Write \\
 & (RPM) &  &  &  & KB \\
\hline
astar & 3619 & 722292 & 9977044 & 8269956 & 10878 \\
gobmk & 3636 & 1263067 & 4075556 & 17123935 & 11622 \\
calculix & 3645 & 697783 & 11007822 & 8025921 & 23544 \\
zeusmp & 3640 & 1271058 & 17166529 & 14100030 & 1052 \\
\hline
\end{tabular}
% END RECEIVE ORGTBL pecs
\begin{comment}
#+ORGTBL: SEND pecs orgtbl-to-latex :splice nil :skip 0
|-----------+-----------+---------+----------+----------+------------|
| Benchmark | Fan Speed |     HT0 |      HT1 |      HT2 | Disk Write |
|           |     (RPM) |         |          |          |         KB |
|-----------+-----------+---------+----------+----------+------------|
| astar     |      3619 |  722292 |  9977044 |  8269956 |      10878 |
| gobmk     |      3636 | 1263067 |  4075556 | 17123935 |      11622 |
| calculix  |      3645 |  697783 | 11007822 |  8025921 |      23544 |
| zeusmp    |      3640 | 1271058 | 17166529 | 14100030 |       1052 |
|-----------+-----------+---------+----------+----------+------------|
\end{comment}

From the listed component data above, we can see that benchmarks
exercise system components at different rates during their execution on
the AMD machine.  In general, benchmark execution affects the full
system, beyond just the processor/cache/memory hierarchy.  The above
table results signify that data move from the processor via the HT2
HyperTransport bus to the disks in a different rate during benchmark
execution.  A similar situation is also observed when the benchmarks are
executed on the Intel machine.

\begin{rviewcomment}
  Sixth Comment/Concern
\end{rviewcomment}
\begin{quote}
  \begin{itshape}
    Last, I am still not clear how these servers exhibit just 60W-70W
    and 45W-70W idle-active power ranges. Are you showing ONLY CPU
    power? Can you please provide more details to how you confirm this
    with data provided by server manufacturers?  Based on what I had
    seen, Sun quick reference reports 450W power consumption for Sun
    Fire 2200 and Dell PowerEdge Power and Performance Data Sheet
    reports Min:138W, Typical:285W, and Max:425W, which are more in line
    with what I expected.
  \end{itshape}
\end{quote}
\rviewresponse

The values reported previously in Figures 7 \& 8 included only CPU
power, as you correctly pointed out. We have redrawn the two figures to
show total system power consumption in this revision.
\clearpage
\bibliographystyle{acmtrans}
\bibliography{../overall}

\end{letter}
\end{document}
% The following comment block is used by the different flavors of EMACS and
% the AUCTEX package to manage multiple documents.  In order for AUCTEX
% to understand you're working with multiple files, you should define
% the TeX-master variable as a file local variable that identifies your
% master document.
%
% Please do not remove.
%%% Local Variables: 
%%% mode: latex
%%% TeX-master: "secondreviewresp.tex"
%%% TeX-PDF-mode: t
%%% End: 



























